% Generated draft: requires mathematical and source review.
\documentclass[11pt]{article}
\usepackage[margin=25mm]{geometry}
\usepackage{fontspec}
\setmainfont{Latin Modern Roman}
\usepackage{amsmath,amssymb,mathtools,unicode-math}
\setmathfont{Latin Modern Math}
\usepackage{enumitem,needspace,xurl,hyperref}
\setlength{\parindent}{0pt}
\setlength{\parskip}{5pt}
\setlength{\emergencystretch}{3em}
\newcommand{\N}{\mathbb{N}}
\newcommand{\Z}{\mathbb{Z}}
\newcommand{\Q}{\mathbb{Q}}
\newcommand{\R}{\mathbb{R}}
\newcommand{\C}{\mathbb{C}}
\newcommand{\F}{\mathbb{F}}
\newcommand{\A}{\mathbb{A}}
\newcommand{\PP}{\mathbb{P}}
\newcommand{\E}{\mathbb{E}}
\newcommand{\eps}{\varepsilon}
\newcommand{\dd}{\,\mathrm{d}}
\begin{document}
\section*{UnsolvedMath ID: 20001051}
\textbf{Source problem code:} catalogue entry
\subsection{Definitions and mathematical statement}
The OCR fragment asks whether $\sum_{i=1}^t\delta^+_{G_i}(v)\geq|V|$ for every vertex forces a rainbow directed cycle, but does not supply the full preceding graph conventions. Under the aid's finite list interpretation, take $n\geq2$ vertices and loopless simple directed layers $G_i$ on the same set. An ordered pair may occur in several layers, giving a list of available labels. The load condition is
\[
 \sum_i d^+_{G_i}(v)\geq n\qquad(v\in V).
\]
A directed cycle is rainbow if its arcs can be assigned distinct available labels. Simplicity within a layer matters: repeated parallel copies of one label must not inflate the load.

For each tail define its active outgoing palette
\[
 P(v)=\{i:d^+_{G_i}(v)>0\},\qquad P(X)=\bigcup_{v\in X}P(v).
\]
The aid first gives a sufficient condition independent of the numerical load: if the palettes satisfy Hall's condition $|P(X)|\geq|X|$ for every $X\subseteq V$, then a rainbow directed cycle exists. A system of distinct label representatives for all tails allows one corresponding outgoing arc at every vertex; the resulting functional digraph has a cycle with distinct labels.

Consequently a counterexample to the load conjecture must have a deficient palette set. For every inclusion-minimal such $X$, the report gives
\[
 |X|\geq3,\qquad |P(X)|=|X|-1.
\]
The size assertion uses $d_i^+(v)\leq n-1$, so the load forces at least two active labels at each vertex. For every chosen root $x\in X$, the full palettes on $X\setminus\{x\}$ admit distinct representatives onto all of $P(X)$. But define the internally realizable palettes on $Y=X\setminus\{x\}$ by
\[
 P_Y(v)=\{i:\exists w\in Y\text{ with }(v,w)\in E(G_i)\}.
\]
These internal palettes must fail Hall's condition; otherwise the same functional selection inside $Y$ would give a forbidden rainbow cycle. Full outgoing labels and labels realizable inside the remaining vertex set cannot be conflated.

The aid further reports the local concentration consequence: every $v\in X$ has some layer $i$ with
\[
 d_i^+(v)\geq\left\lceil\frac n{|X|-1}\right\rceil.
\]
This is a necessary property of a hypothetical counterexample, not a construction of one.

A separate sufficient balanced-support criterion is obtained by writing
\[
 b=\min_v|P(v)|>0,\qquad
 D=\max_i\#\{v:i\in P(v)\}.
\]
If $b\geq D$, incidence counting implies Hall's condition and therefore a rainbow directed cycle. The positive-palette assumption holds under the stated load condition. These palette certificates and obstructions do not resolve all remaining cases of the original conjecture.
\subsection{Short English statement}
Distinct outgoing labels for all vertices immediately force a rainbow cycle. The aid uses this to constrain every possible counterexample: a minimal deficient palette set has deficiency exactly one, and its internal palettes fail even though its full palettes can be matched after deleting a root. A balanced label-support condition gives a separate sufficient case.
\subsection{Sources}
\begin{itemize}
\item Local UnsolvedMath registry, permanent ID $20001051$, source code \texttt{catalogue entry}, supplied OCR fragment and research summary.
\item American Institute of Mathematics, \emph{The Caccetta--Haggkvist conjecture}, source-item label 6.1, supplied URL \url{https://aimath.org/WWN/caccetta/caccetta.pdf}; canonical local record \texttt{aim-combinatorics-notes.json}, \texttt{notes[175]}. No PDF was retrieved.
\item The original literature field is empty. No full authoritative bibliographic reference for the reconstructed list conventions, palette-transversal lemma or counterexample constraints is supplied.
\item Research attempt 1: partial result; problem open at the recorded run, registry status partially solved; supplied rigor label ``proved'', novelty confidence ``low''. The machine-generated aid is not peer reviewed and requires independent verification before citation.
\end{itemize}
\end{document}
